\section{Политический эксперимент}

В эксперименте строго разделяются истинное локальное состояние, наблюдения, информационное
состояние, правило политики и критерий качества. Истинное состояние \(q_t\) порождается
HANK/SSJ-динамикой. Центральный банк наблюдает только \(o_t\), строит информационное состояние
\(s_t^\phi=\phi(\mathcal I_t)\), выбирает ставку по правилу \(\mu_\theta(s_t^\phi)\), а потери
оцениваются по истинным реализованным переменным из \(q_t\).

Правило ставки строится по выбранному информационному состоянию:
\[
i_t=\mu_\theta(s_t^\phi).
\]
В основной спецификации используется интерпретируемое линейное правило
\[
i_t
=
\rho_i i_{t-1}
+
\theta^\top s_t^\phi.
\]
Для каждого информационного состояния коэффициенты подбираются заново:
\[
\theta_\phi^*
=
\arg\min_{\theta\in\Theta}
J(\theta,\phi).
\]
Это важно: сравнение информационных состояний не должно смешиваться с тем, что одно правило было настроено лучше другого.

В финальной спецификации используется единый бюджет оптимизации для всех информационных
состояний. Случайные кандидаты и начальные точки, построенные по локально SSJ-оптимальной ставке,
используются только как стартовые точки. После этого для каждого информационного состояния
запускается одна и та же многозапусковая непрерывная оптимизация с одинаковым числом стартов,
одинаковым числом итераций, одинаковыми методами, одинаковыми границами коэффициентов и одной и
той же сеткой ridge-регуляризации. Поэтому различия в качестве правил не могут быть объяснены
тем, что более богатое информационное состояние получило больший бюджет подбора.

Функция потерь является квадратичным стабилизационным критерием:
\[
L_t
=
\lambda_\pi \pi_t^2
+
\lambda_y y_t^2
+
\lambda_c C_t^2
+
\lambda_i(i_t-i_{t-1})^2.
\]
Накопленные потери:
\[
J(\theta,\phi)
=
\mathbb E
\sum_{t=0}^{T}
\beta^t L_t.
\]
Эта величина не называется welfare-loss; она используется как прозрачный критерий стабилизационного качества правила.

Центральная проверяемая гипотеза формулируется как равенство потерь у правила по фильтрованным
агрегатам и у правила по фильтрованному распределительному состоянию:
\[
H_0:
\quad
J(s^{filt-agg})=J(s^{filt-dist}).
\]
Содержательная альтернатива:
\[
H_1:
\quad
J(s^{filt-dist})<J(s^{filt-agg}),
\]
но не как универсальное утверждение. Ожидаемый эффект должен проявляться именно в средах, где
распределительные статистики помогают оценить будущую трансмиссию ставки.

Главная метрика работы:
\[
MVOI_{dist}
=
J(s^{filt-agg})-J(s^{filt-dist}).
\]
Дополнительно считается общий выигрыш относительно агрегатных наблюдений:
\[
VOI_{dist}
=
J(s^{obs-agg})-J(s^{filt-dist}),
\]
и доля сокращения измеренного разрыва до ориентира полной информации:
\[
ShareClosed(\phi)
=
\frac{J(s^{obs-agg})-J(s^\phi)}
{J(s^{obs-agg})-J(s^{full})}.
\]

Основные эксперименты строятся на траекториях, порождённых HANK-моделью и методом последовательностей.
В основной спецификации фильтрованные информационные состояния строятся совместным линейным
фильтром Калмана:
\[
q_{t+1}=Aq_t+\varepsilon_{t+1},
\qquad
o_t=M_\phi q_t+\nu_t.
\]
Матрица наблюдения \(M_\phi\) задаёт информационную структуру: агрегатные сигналы, агрегатные и
распределительные сигналы или полную информацию. При этом матрица перехода \(A\) и ковариация
шоков оцениваются один раз по HANK/SSJ-траекториям и остаются общими для всех информационных
состояний. Это отделяет вопрос о ценности информации от вопроса о том, была ли для одного набора
переменных задана другая динамика.

Скалярная фильтрация отдельных рядов не используется как основная спецификация. Она остаётся
только технической проверкой, показывающей, насколько результаты чувствительны к упрощённому
способу построения фильтрованных переменных.

Перед интерпретацией результатов проверяется локальная линейная аппроксимация HANK/SSJ. Для
денежного шока нелинейный переходный отклик сравнивается с экспортированным SSJ-якобианом. Для
неполитических шоков, по которым текущий архив не хранит готовые SSJ-матрицы, проверяется
конечная разностная локальная линейность переходного решателя. Эта проверка задаёт область, в
которой последующие контрфактические расчёты можно читать как локальные около стационарного
состояния.

Чтобы проверить, что вывод не является артефактом одного класса правил, дополнительно оцениваются
альтернативные спецификации правила ставки: правило типа Тейлора, правило Тейлора с потреблением,
правило Тейлора с распределительными показателями, правило с ограничениями на знаки, а также
регуляризованные линейные правила. Малое нелинейное правило используется только как дополнительная
проверка и не входит в центральный вклад работы.

Дополнительно проводится усиленная проверка с обратной связью. В базовой локальной проекции правило
выбирает траекторию ставки по информационным признакам исходной HANK/SSJ-траектории, после чего
агрегатные переменные пересчитываются через SSJ-отклик. В проверке с обратной связью траектория
ставки, контрфактическое состояние, шумные наблюдения и фильтрованные признаки пересчитываются
итеративно. Для распределительных статистик используются прямые HANK/SSJ-отклики на траекторию
ставки. Правила при этом не переобучаются: используются коэффициенты, выбранные в основном
эксперименте. Поэтому этот блок является проверкой устойчивости уже найденных правил, а не новой
оптимизацией политики.

Наконец, для той же совместной линейной state-space задачи строится ориентир Riccati. В этой
проверке используется локальная система
\[
q_{t+1}=Aq_t+Bi_t+\varepsilon_{t+1},
\qquad
o_t=M_\phi q_t+\nu_t,
\]
где \(A,Q,M_\phi,R_\phi\) берутся из совместной спецификации фильтра Калмана, а \(B\) оценивается
по HANK/SSJ-траекториям как локальная связь между ставкой и следующим состоянием. Для полной
информации решается конечногоризонтная LQR-задача, а для неполной информации -- соответствующая
LQG-задача с теми же матрицами наблюдения, что и в основном эксперименте. Этот блок не заменяет
основные простые правила, а отвечает на отдельный вопрос: насколько далеко ограниченные правила
находятся от оптимального линейного регулятора в той же локальной state-space задаче.

\subsection{Почему это не просто добавление переменных в регрессию}

Главное сравнение устроено не как произвольное расширение регрессора. Во-первых, для всех
информационных состояний используется один и тот же класс правил -- линейное правило ставки с
заново подобранными коэффициентами. Во-вторых, параметры каждого правила выбираются только на
валидационных траекториях, а итоговые потери считаются на независимых тестовых траекториях.
В-третьих, сравнение является парным: разные информационные состояния оцениваются на одинаковых
реализациях HANK/SSJ-шоков и одинаковых шумах наблюдений. Поэтому разность потерь отражает
ценность информационного состояния, а не различие в классе правил или удачный набор тестовых
траекторий. Наконец, проверки с перемешанными и искусственными распределительными рядами
показывают, исчезает ли эффект, когда распределительный сигнал перестаёт быть связан с
HANK/SSJ-динамикой.
